\section{Discussion}
\label{sec:discussion}

\subsection{Summary of Main Findings}
\label{sec:discussion-summary}

Pooled across the seventeen studies contributing clinical-success data, patients treated with bacteriophage therapy for MDR, XDR, or PDR \textit{P.\ aeruginosa} infections achieved clinical success in 76.9\% of cases (95\% CI 64.4--86.0\%; \Cref{fig:forest-success-overall}), directionally in line with earlier, non-stratified syntheses of this literature \parencite{ElHaddad2019_eskape,Uyttebroek2022_lancetid,Liu2025_ijaa}. That figure describes a compassionate-use population offered phage therapy after antibiotics had already failed, not a treatment effect measured against a comparator, and it is reported here as this review's secondary, contextual estimand rather than its headline finding (Section~\ref{sec:results-synthesis}). The headline finding is a partial success on the stratification this review was built to test: of the 37 outcome-stratum cells defined by resistance category, route, and modality, 16 met the pre-specified pooling threshold, including, for the first time, a clean resistance-specific stratum --- MDR, pooled across all four outcomes (11--14 studies, 35--44 patients depending on outcome) --- alongside a heterogeneous multi-route ``other'' bucket, safety and mortality within the residual not-classifiable resistance bucket, and, newly, safety and mortality in the inhaled/nebulized route (\Cref{tab:pooled-estimates,tab:not-pooled}). That MDR estimate answers part of the question this review set out to ask. It does not complete it: XDR (4 studies, 10 patients) and PDR (3--4 studies, 4--6 patients) remain below the patient threshold for every outcome, so the MDR-\emph{versus}-XDR comparison that motivated the stratification in the first place still cannot be made, and the entire monotherapy-versus-combination comparison still fell short entirely (0 antibiotic-monotherapy arms). This review's contribution is documenting exactly which parts of that stratification the evidence can and cannot yet support -- neither forcing an estimate it cannot bear nor dismissing the one clean success as noise.

\subsection{Comparison with Prior Syntheses}
\label{sec:discussion-comparison}

\textcite{ElHaddad2019_eskape}'s review of ESKAPE-pathogen phage therapy, \textcite{Uyttebroek2022_lancetid}'s larger \textit{Lancet Infectious Diseases} synthesis, and \textcite{Liu2025_ijaa}'s 130-study individual-patient-data meta-analysis each report favorable pooled outcomes for phage therapy broadly, consistent in direction with the success proportions found here, though none restricts its evidence base to \textit{P.\ aeruginosa} specifically or reports a directly comparable resistance- or pathogen-specific figure against which 76.9\% (overall) or 74.3\% (MDR-specific) could be checked numerically. What none of the three does is stratify by resistance category or route. This review can now do so partially: the MDR-specific pooled estimate (74.3\% clinical success, 95\% CI 55.4--87.0\%) is, to our knowledge, the first resistance-stratified pooled proportion for phage therapy in \textit{P.\ aeruginosa} specifically to clear a pre-specified evidentiary threshold. It remains an incomplete answer to the question that motivated the stratification --- a clinician still cannot compare it against an XDR- or PDR-specific figure, since neither clears the same threshold, and combination-versus-monotherapy remains entirely unanswerable (zero antibiotic-monotherapy arms) --- but identifying, cell by cell, which of those comparisons the published record can and cannot yet support is where this review adds something those three broader syntheses did not attempt.

\subsection{Safety Findings and Small-Study Effects}
\label{sec:discussion-safety}

The pooled safety estimate illustrates a problem with this literature more broadly, not only with this review's own numbers. The GLMM estimate (19.4\%, 95\% CI 10.0--34.4\%; Section~\ref{sec:results-synthesis}) is bracketed by its sensitivity models (Freeman-Tukey 8.2\%, logit-DerSimonian-Laird and fixed-effect both 26.9\%), and leave-one-out does not flag the \emph{overall} safety estimate as sensitive to any single study's removal. The MDR-specific safety and mortality estimates, the route-``other'' safety estimate, and the not-classifiable resistance stratum's mortality estimate remain sensitive, however, collapsing to a value indistinguishable from zero if \citeauthor{Pirnay2024_natmicrobiol} or \citeauthor{Jault2019_phagoburn} is removed. The MDR resistance stratum's pooled estimate therefore still leans on \citeauthor{Pirnay2024_natmicrobiol}'s cohort (23 of the 35 MDR patients in the clinical-success pool, 66\%; a smaller 56\% of the larger 41-patient MDR eradication pool), though the seven MDR arms added by the native Scopus search dilute that single-source dependence rather than deepen it, corroborated rather than counterbalanced by a set of much smaller studies.

Egger's test, applied to the outcome-overall cells with at least ten studies as pre-specified, now covers four outcomes. Three are non-significant --- clinical success ($p=0.26$), safety ($p=0.93$), and eradication ($p=0.37$; newly eligible now that eradication clears the ten-study minimum) --- but mortality is \textbf{significant} ($p=0.02$; \Cref{fig:funnel-eyeball}), and more so within the MDR stratum ($p=0.005$). The enlarged corpus thus reverses the earlier finding of no significant small-study asymmetry for at least one outcome: two small single-patient fatal case reports added by the native Scopus search (Lev\^eque et al.\ 2023; Duplessis et al.\ 2018) introduced exactly the asymmetry a mortality funnel is most prone to, with small studies reporting deaths clustering at the low-precision end. We read this as a signal that the pooled mortality figure (10.7\%) may if anything under-state true mortality, not as grounds to dismiss it --- it is one significant test among several non-significant ones in a still-small corpus. The other three outcomes' non-significant tests, meanwhile, remain weak reassurance rather than a clean acquittal: this corpus of 28 study-arms remains far smaller than any of the three prior syntheses discussed above, single-patient case reports still make up the largest design category (Section~\ref{sec:results-characteristics}), and the classification-sensitivity check (Section~\ref{sec:results-robustness}) shows the MDR finding depends on author-reported, not independently re-verified, resistance labels. Where a formal test does detect asymmetry, as for mortality, it should be taken seriously; where it does not, in a corpus this small and this dependent on one or two dominant sources, its silence is weak evidence of absence. We read the combined picture as narrowing rather than eliminating the small-study-effects concern a referee reviewing case-report-heavy evidence syntheses would be right to raise, and as reason to read the broader literature's high-success narrative, including the point estimates reported by \textcite{ElHaddad2019_eskape}, \textcite{Uyttebroek2022_lancetid}, and \textcite{Liu2025_ijaa}, with real caution.

\subsection{Limitations}
\label{sec:discussion-limitations}

Four limitations bear most directly on how much weight these estimates can carry. First, screening and extraction were conducted by a single reviewer, with a second reviewer's critique applied sequentially rather than concurrently (Sections~\ref{sec:methods-selection}--\ref{sec:methods-extraction}). That process caught several substantive errors before they entered, or after they had briefly entered, the analytic dataset --- a study initially catalogued as the closest candidate antibiotic-monotherapy comparator in this search, later confirmed on full text to contain zero \textit{P.\ aeruginosa} cases; a duplicate patient reported under two citations; and, most recently, two patients within \citeauthor{Pirnay2024_natmicrobiol}'s cohort whose resistance profile, once checked against per-patient data, turned out not to meet this review's own MDR/XDR/PDR eligibility criterion (Section~\ref{sec:results-characteristics}); and, from the native Scopus search, a single XDR case (Van Nieuwenhuyse et al., 2022) recognized as already counted within \citeauthor{Pirnay2024_natmicrobiol}'s Belgian-consortium cohort and excluded to prevent double-counting --- evidence the sequential-critique and re-extraction process catches real errors, and that the consortium double-counting risk is real and was actively screened for, but not a substitute for the dual-independent process a registered protocol would specify; a formal re-screen could still change which studies are included. Second, risk-of-bias ratings, now complete for all 29 study-arms (Section~\ref{sec:results-rob}), surface a genuine concern rather than resolve one. The trial that anchors this review's phage-monotherapy-adjacent evidence throughout --- PhagoBurn \parencite{Jault2019_phagoburn}, the only source with a phage-alone treatment arm compared against a non-phage standard-of-care --- rates \textbf{High risk of bias} on RoB2, driven by unmasked clinicians and by the trial stopping early for insufficient efficacy, both recognized bias-introducing deviations. This does not overturn anything already reported (PhagoBurn's clinical-success data were already excluded from pooling as a time-to-event, non-binary endpoint), but it does mean the one study in this corpus that most directly speaks to phage-versus-standard-of-care performance is also the one whose result should be trusted least on methodological grounds --- a genuinely unfavorable combination for a literature already short on comparative evidence. \citeauthor{Pirnay2024_natmicrobiol}'s cohort, by contrast, rates Low risk on the JBI checklist (explicitly consecutive, complete, transparent per-patient reporting), which if anything strengthens rather than weakens confidence in its now-dominant role in the MDR-specific estimates (Section~\ref{sec:discussion-safety}) --- the concern with that source is its weight in the pooled MDR estimate, not its individual quality. Third, the MDR resistance stratum --- this review's headline new finding --- still leans on a single source: \citeauthor{Pirnay2024_natmicrobiol} supplies 23 of the 35 MDR patients in the clinical-success pool (66\%; the several MDR arms added by the native Scopus search dilute this to 56\% in the larger 41-patient MDR eradication pool), and a classification-sensitivity check shows the stratum would fail the pooling threshold entirely (dropping to 3 studies, 4 patients) if restricted to independently-verified rather than author-reported resistance labels (Section~\ref{sec:results-robustness}). The MDR estimate should therefore be read as characterizing one large cohort's outcomes, corroborated in direction by a set of much smaller studies that now dilute --- but do not overturn --- its single-source dependence, not as an even-weighted synthesis across independent sources --- and XDR and PDR remain too sparse (10 and 4--6 patients respectively, depending on outcome) to pool at all, so the resistance-class comparison this review was built to answer is still only partially available. Fourth, the database coverage, though now materially improved, is not exhaustive. A native Scopus search (Section~\ref{sec:methods-search}) was added as a third core database alongside PubMed/MEDLINE and ClinicalTrials.gov and roughly doubled the pooling-eligible corpus (from 13 to 23 studies), adding ten studies the earlier PubMed-only rounds had missed and correcting one prior exclusion --- which both confirms the earlier coverage gap was real and material and largely closes it, since three searched databases meet a recognized minimum for a defensible PRISMA search and Scopus indexes much of EMBASE's content. EMBASE and Web of Science remain unsearched, but this is now a minor residual gap rather than a headline limitation. That the pooled eradication estimate has shifted by more than ten percentage points as studies were added across search rounds (standing at 55.5\% in the current corpus; \Cref{tab:pooled-estimates}) is nonetheless a direct reminder that the corpus is still small enough for a handful of additional records to move a headline figure, and completing a native EMBASE and Web of Science search would further stabilize the current estimates. Two data-handling limitations attach to the newly added studies: Chan et al.'s adverse-event data are reported only at the nine-patient cohort level and could not be attributed to the MDR-versus-PDR sub-groups, so they are recorded but not carried as a usable safety numerator (a documented data loss that, among other things, left the inhaled/nebulized route just short of the safety-pooling threshold); and Chan et al.'s compassionate cohort may share patients with the Yale CYPHY randomized trial, which is one reason CYPHY is kept out of pooling. Two further limitations are disclosed for completeness: no PROSPERO registration before extraction began (Section~\ref{sec:methods-protocol}), and a corpus, at 28 pooling-eligible arms and 106 patients, still too small to support most of the stratification it was designed to test --- the limitation this review treats as its central finding, not a footnote.

\subsection{Implications for Practice and Research}
\label{sec:discussion-implications}

For clinicians, the pooled proportions reported here should not inform a treatment decision: they summarize a referral-selected, salvage-therapy population, the resistance-class stratification that succeeded (MDR) rests overwhelmingly on a single cohort, and route- and modality-specific comparisons remain unanswerable. Most strata in this evidence base are expected to warrant a Very Low GRADE certainty rating once that assessment is completed (Section~\ref{sec:methods-rob}). For the field, this review's own experience points to three concrete needs: standardized outcome definitions across primary studies, since extracting ``clinical success'' required accepting each study's own definition because no consistent one exists (Section~\ref{sec:methods-eligibility}); comparative studies with a genuine antibiotic-monotherapy arm for \textit{P.\ aeruginosa} specifically, since despite a targeted search essentially none exist, and additional single-arm case reporting will not close that gap; and retrospective PROSPERO registration of this review disclosing the sequence in Section~\ref{sec:methods-protocol}, with prospective registration for any successor synthesis in this space.

\subsection{Conclusion}
\label{sec:discussion-conclusion}

What this review can defensibly claim is narrower than a single pooled percentage, but broader than a review that could not stratify by resistance category at all: bacteriophage therapy, used adjunctively in a small, referral-selected group of patients with MDR, XDR, or PDR \textit{P.\ aeruginosa} infections refractory to antibiotics, was associated with a high proportion of reported clinical success, including now a resistance-specific (MDR) estimate that, to our knowledge, is the first of its kind for this pathogen to clear a pre-specified evidentiary threshold --- though that estimate is itself dominated by one large cohort, not an even-weighted synthesis. What remains unknown is most, but no longer all, of what the stratification was designed to resolve: whether MDR, XDR, and PDR isolates respond differently is still unanswerable, since neither XDR nor PDR can yet be pooled at all, and whether route of administration or combination therapy matters remains equally open. The single most useful next step for this field is not another synthesis of the existing case-report literature, but two complementary efforts: independent replication of the MDR resistance-stratified estimate reported here in additional cohorts, so it no longer rests on one source, and a comparative study, ideally randomized, that pairs a genuine antibiotic-monotherapy arm with a phage-antibiotic combination arm in an MDR-, XDR-, or PDR-\textit{Pseudomonas}-specific population. Until both exist, the question this review set out to answer will remain only partially resolved.
